\chapter{Integral Dependence and Valuations}

\section*{INTEGRAL DEPENDENCE}

\begin{proposition}{5.1}
	The following are equivalent:
	\begin{enumerate}[label=\textup{(\roman*)}]
		\item $x\in B$ is integral over $A$;
		\item $A[x]$ is a finitely generated $A$-module;
		\item $A[x]$ is contained in a subring $C$ of $B$ such that $C$ is a finitely generated $A$-module;
		\item There exist a faithful $A[x]$-module $M$ which is finitely generated as an $A$-module.
	\end{enumerate}
\end{proposition}

\begin{definition}
	The set $\overline{A}$ of elements of $B$ which are integral over $A$ is called the \textcolor{orange}{integral closure} of $A$ in $B$. If $\overline{A}=A$, then $A$ is said to be \textcolor{orange}{integrally closed in $B$}. If $\overline{A}=B$, the ring $B$ is said to be \textcolor{orange}{integral over $A$}.
\end{definition}

\section*{THE GOING-UP THEOREM}
\begin{proposition}{5.7}
	Let $A\subseteq B$ be integral domains, $B$ integral over $A$. Then $B$ is a field iff $A$ is a field.
\end{proposition}

\begin{corollary}{5.9}
	Let $A\subseteq B$ be rings, $B$ integral over $A$. Let $\mathfrak{q}_1$, $\mathfrak{q}_2$ be prime ideals of $B$ such that $\mathfrak{q}_1\subseteq\mathfrak{q}_2$ and $\mathfrak{q}_1^c=\mathfrak{q}_2^c=\mathfrak{p}$. Then $\mathfrak{q}_1=\mathfrak{q}_2$.
\end{corollary}

\begin{theorem}{5.10}[\textcolor{red}{Lying over property}]
	Let $A\subseteq B$ be rings, $B$ integral over $A$, and let $\mathfrak{p}$ be a prime ideal of $A$. Then there exists a prime ideal $\mathfrak{q}$ of $B$ such that $\mathfrak{q}\cap A=\mathfrak{p}$.
\end{theorem}

\begin{theorem}{5.11}[Going-up theorem]
	Let $A\subseteq B$ be rings, $B$ integral over $A$. Let $\mathfrak{p}_1\subseteq\cdots\subseteq\mathfrak{p}_n$ be a chain of prime ideals of $A$ and $\mathfrak{q}_1\subseteq\cdots\subseteq\mathfrak{q}_m$ $(m<n)$ a chain of prime ideals of $B$ such that $\mathfrak{q}_i\cap A=\mathfrak{p}_i$ for $1\leqslant i\leqslant m$. Then the chain $\mathfrak{q}_1\subseteq\cdots\subseteq\mathfrak{q}_m$ can be extended to a chain $\mathfrak{q}_1\subseteq\cdots\subseteq\mathfrak{q}_n$ such that $\mathfrak{q}_i\cap A=\mathfrak{p}_i$ for $1\leqslant i\leqslant n$.
\end{theorem}

\section*{THE GOING-DOWN THEOREM}

\begin{proposition}{5.12}
	Let $A\subseteq$ be rings, $C$ the integral closure of $A$ in $B$. Let $S$ be a multiplicatively closed subset of $A$. Then $S^{-1}C$ is the integral closure of $S^{-1}A$ in $S^{-1}B$.
\end{proposition}

An integral domain is said to be \textcolor{orange}{integrally closed} (without qualification) if it is integrally closed in its field of fractions.
	
\begin{proposition}{5.13}
	Let $A$ be an integral domain. Then the following are equivalent:
	\begin{enumerate}[label=\textup{(\roman*)}]
		\item $A$ is integrally closed;
		\item $A_\mathfrak{p}$ is integrally closed, for each prime ideal $\mathfrak{p}$;
		\item $A_\mathfrak{m}$ is integrally closed, for each maximal ideal $\mathfrak{m}$.
	\end{enumerate}
\end{proposition}

\begin{theorem}{5.16}[Going-down theorem]
	Let $A\subseteq B$ be integral domains, $A$ integrally closed, $B$ integral over $A$. Let $\mathfrak{p}_1\supseteq\cdots\supseteq\mathfrak{p}_n$ be a chain of prime ideals of $A$ and $\mathfrak{q}_1\supseteq\cdots\supseteq\mathfrak{q}_m$ $(m<n)$ a chain of prime ideals of $B$ such that $\mathfrak{q}_i\cap A=\mathfrak{p}_i$ for $1\leqslant i\leqslant m$. Then the chain $\mathfrak{q}_1\supseteq\cdots\supseteq\mathfrak{q}_m$ can be extended to a chain $\mathfrak{q}_1\supseteq\cdots\supseteq\mathfrak{q}_n$ such that $\mathfrak{q}_i\cap A=\mathfrak{p}_i$ for $1\leqslant i\leqslant n$.
\end{theorem}

\section*{VALUATION RINGS}

\begin{definition}
	Let $B$ be an integral domain. $B$ is a \textcolor{orange}{valuation ring} of $\operatorname{Frac}(B)$ if, for each $x\neq0$, either $x\in B$ or $x^{-1}\in B$ (or both).
\end{definition}

\begin{proposition}{5.18}
	Let $B$ be a valuation ring, $K=\operatorname{Frac}(B)$ its field of fractions. Then
	\begin{enumerate}[label=\textup{(\roman*)},ref=\thepropositiontemp(\roman*)]\crefalias{enumi}{proposition}
		\item\label{5.18.1} $B$ is a local ring;
		\item\label{5.18.2} If $B^\prime$ is a ring such that $B\subseteq B^\prime\subseteq K$, then $B^\prime$ is a valuation ring of $K$;
		\item\label{5.18.3} $B$ is integrally closed.
	\end{enumerate}
\end{proposition}

\begin{theorem}{5.21}
	Let $K$ be a field, $\Omega$ an algebraically closed field. Let $\Sigma$ be the set of all pairs $(A,f)$, where $A$ is a subring of $K$ and $f:A\rightarrow\Omega$ a homomorphism. Define a partial order of $\Sigma$ to be:
	\begin{equation*}
		(A,f)\leqslant(B,g)\Leftrightarrow A\subseteq B \text{ and } g|_A=f.
	\end{equation*}
	Then the maximal element of $\Sigma$ is a valuation ring of $K$.
\end{theorem}

\section*{DEFINITIONS AND PROPOSITIONS FROM EXERCISES}

\begin{definition}
	A ring homomorphism $f: A\rightarrow B$ is said to have the \textcolor{orange}{going-up property} if for $\mathfrak{p}_1\subseteq\mathfrak{p}_2$ and $\mathfrak{p}_1=f^{-1}(\mathfrak{q}_1)$, there exist a $\mathfrak{q}_2$ such that $\mathfrak{q}_1\subseteq\mathfrak{q}_2$ and $f^{-1}(\mathfrak{q}_2)=\mathfrak{p}_2$. Reversing the inclusions of prime ideals we have the \textcolor{orange}{going-down property}.
\end{definition}

\begin{proposition*}[\textcolor{red}{Noether normalization lemma}]
	Let $k$ be a field and let $A\neq0$ be a finitely generated $k$-algebra. Then there exist elements $y_1,\cdots,y_r\in A$ which are algebraically independent over $k$ and such that $A$ is integral over $k[y_1,\cdots,y_r]$.
\end{proposition*}

\begin{definition}
	A ring $A$ satisfied the following equivalent property is a \textcolor{orange}{Jacobson ring}.
	\begin{enumerate}[label=\textup{(\roman*)}]
		\item Every prime ideal in $A$ is an intersection of maximal ideals.
		\item In every homomorphic image of $A$ the nilradical is equal to the Jacobson radical.
		\item Every prime ideal in $A$ which is not maximal is equal to the intersection of the prime ideals which contain it strictly.
	\end{enumerate}
\end{definition}

\begin{definition}
	Let $A$ be a valuation ring of a field $K$. The group $U$ of units of $A$ is a subgroup of the multiplicative group $K^*$ of $K$. Let $\Gamma=K^*/U$. If $\xi,\iota\in\Gamma$ are represented by $x,y\in K$, define $\xi\geqslant\iota$ to mean $xy^{-1}\in A$. Then $\Gamma$ is a totally ordered abelian group. It is called the \textcolor{orange}{value group} of $A$.
	
	Conversely, let $\Gamma$ be a totally ordered abelian group. A \textcolor{orange}{valuation} of $K$ with values in $\Gamma$ is a mapping $v:K^*\rightarrow\Gamma$ such that
	\begin{enumerate}[label=\textup{(\arabic*)}]
		\item $v(xy)=v(x)+v(y)$,
		\item $v(x+y)\geqslant\operatorname{min}(v(x),v(y))$,
	\end{enumerate}
	for all $x,y\in K^*$. Then the set of elements $x\in K^*$ such that $v(x)\geqslant0$ is a valuation ring of $K$. This ring is called the \textcolor{orange}{valuation ring} of $v$, and the subgroup $v(K^*)$ of $\Gamma$ is the \textcolor{orange}{value group} of $v$. Thus the concepts of valuation ring and valuation are equivalent.
\end{definition}